\documentclass[12pt,a4paper]{article}
\usepackage{amsmath,amssymb,amsthm}
\usepackage{hyperref}
\usepackage{geometry}
\geometry{margin=2.5cm}
\usepackage{enumitem}
\usepackage{booktabs}

\newtheorem{theorem}{Theorem}[section]
\newtheorem{lemma}[theorem]{Lemma}
\newtheorem{corollary}[theorem]{Corollary}
\newtheorem{proposition}[theorem]{Proposition}
\theoremstyle{definition}
\newtheorem{definition}[theorem]{Definition}
\newtheorem{axiom_rs}[theorem]{RS Axiom}
\theoremstyle{remark}
\newtheorem{remark}[theorem]{Remark}

\title{The Neutron Star Maximum Mass\\
and the Tolman-Oppenheimer-Volkoff Limit\\
from Recognition Science}

\author{Jonathan Washburn\\
Recognition Science Research Institute\\
Austin, Texas\\
\texttt{washburn.jonathan@gmail.com}}

\date{March 2026}

\begin{document}
\maketitle

\begin{abstract}
The Tolman-Oppenheimer-Volkoff (TOV) equation governs the hydrostatic
equilibrium of neutron stars in General Relativity.  Combined with a
nuclear equation of state (EOS), it predicts a maximum neutron star mass
$M_\mathrm{TOV}$.  Observations of PSR~J0740+6620 ($2.08 \pm 0.07\,
M_\odot$) and PSR~J0952-0607 ($2.35 \pm 0.17\,M_\odot$) constrain
$M_\mathrm{TOV} \gtrsim 2.0\,M_\odot$.  We derive the TOV equation from
the Recognition Science (RS) framework, in which GR emerges from the
RS action principle (Lean: \texttt{Relativity.Dynamics.EFEEmergence}),
and the neutron EOS is derived from the $\varphi$-ladder mass and the
degenerate neutron pressure.  The RS prediction for the maximum mass is
$M_\mathrm{TOV}^\mathrm{RS} \approx 2.0$--$2.5\,M_\odot$ (depending on
the neutron-neutron interaction rung correction), consistent with
observations.  This paper complements
\texttt{RS\_Chandrasekhar\_Mass.tex} (white dwarfs, $M_\mathrm{Ch} =
1.44\,M_\odot$): the TOV limit is the neutron-star analogue, replacing
electron degeneracy pressure with neutron degeneracy pressure and
requiring the full GR treatment absent in the Chandrasekhar case.
\end{abstract}

%%============================================================
\section{Introduction}
\label{sec:intro}
%%============================================================

The maximum mass of a neutron star is one of the most important
observables in nuclear astrophysics.  Unlike the Chandrasekhar limit
for white dwarfs ($M_\mathrm{Ch} = 1.44\,M_\odot$, where Newtonian
gravity + special relativity suffice), the neutron star maximum mass
requires the full machinery of General Relativity because neutron star
central densities exceed $\sim 2\rho_\mathrm{nuc}$ ($\rho_\mathrm{nuc}
\approx 2.3 \times 10^{14}$ g/cm$^3$).

The Tolman-Oppenheimer-Volkoff (TOV)
equation~\cite{TOV1939,Oppenheimer1939} is the GR generalization of
hydrostatic equilibrium:
\begin{equation}
  \frac{dP}{dr} = -\frac{[\varepsilon(r) + P(r)][M(r) + 4\pi r^3 P(r)/c^2]}
  {r^2[1 - 2GM(r)/(c^2 r)]},
  \label{eq:tov}
\end{equation}
with the mass equation $dM/dr = 4\pi r^2 \varepsilon(r)/c^2$, where
$\varepsilon = \rho c^2$ is the energy density and $P = P(\varepsilon)$
is the EOS.

%%============================================================
\section{Recognition Science Framework}
\label{sec:rs}
%%============================================================

\subsection{GR Hydrostatic Equilibrium from RS}

The TOV equation is the static, spherically symmetric limit of the
GR Euler equation for a perfect fluid.  In the RS framework:
\begin{enumerate}
  \item The EFE emerge from RS action stationarity
        (Lean: \texttt{Relativity.Dynamics.EFEEmergence.efe\_from\_mp}).
  \item The Schwarzschild metric is the unique static spherically
        symmetric solution (Lean: \texttt{Relativity.StaticSpherical}).
  \item For a perfect fluid, the stress-energy tensor is
        $T^{\mu\nu} = (\varepsilon + P)u^\mu u^\nu + Pg^{\mu\nu}$, and
        its conservation $\nabla_\mu T^{\mu\nu} = 0$ in the static
        Schwarzschild interior gives the TOV equation.
\end{enumerate}

\begin{theorem}[TOV equation from RS]\label{thm:tov}
For a static, spherically symmetric, perfect-fluid configuration in
the RS-derived Schwarzschild metric, conservation of the stress-energy
tensor $\nabla_\mu T^{\mu\nu} = 0$ yields
\begin{equation}
  \frac{dP}{dr} = -\frac{G(\varepsilon + P)(M + 4\pi r^3 P/c^2)}
  {c^2 r^2(1 - 2GM/c^2 r)}.
\end{equation}
From Lean modules: \texttt{Relativity.Dynamics.EFEEmergence} $+$
\texttt{Relativity.StaticSpherical}.
\end{theorem}

Note the two GR corrections vs.\ Newtonian hydrostatics:
\begin{itemize}
  \item $\varepsilon + P$ replaces $\varepsilon$ (pressure contributes
        to gravitational source — purely relativistic).
  \item $M + 4\pi r^3 P/c^2$ replaces $M$ and $(1 - 2GM/c^2 r)^{-1}$
        in the denominator.
\end{itemize}
At low densities, $P \ll \varepsilon$ and $r \gg r_s$, and the TOV
equation reduces to Newtonian hydrostatics.

\subsection{Neutron EOS from the $\varphi$-Ladder}

The neutron star EOS in RS is derived from the $\varphi$-ladder mass
of the neutron and the degenerate neutron pressure.

\begin{definition}[Neutron mass on the $\varphi$-ladder]
The neutron mass in RS is
\begin{equation}
  m_n = m_\mathrm{struct} \cdot \varphi^{r_n - 8 + \mathrm{gap}(Z_n)},
\end{equation}
where $r_n$ is the neutron rung, $Z_n$ is the neutron charge-band
parameter (for the neutron, $Z_n = 0$ for the isospin-neutral component),
and $m_\mathrm{struct}$ is the sector yardstick.  Numerically,
$m_n = 939.565$ MeV/$c^2$ (Lean: \texttt{Masses.MassLaw.predict\_mass}).
\end{definition}

For a free, non-interacting degenerate neutron gas (valid near nuclear
saturation density), the EOS is:
\begin{align}
  \varepsilon &= \frac{m_n^4 c^5}{3\pi^2\hbar^3}
  \left[x\sqrt{1+x^2}(2x^2/3 - 1) + \ln(x + \sqrt{1+x^2})\right], \\
  P &= \frac{m_n^4 c^5}{3\pi^2\hbar^3}
  \left[x\sqrt{1+x^2}(2x^2/3 - 1) - \frac{1}{3}x\sqrt{1+x^2}
  + \tfrac{1}{3}\ln(x+\sqrt{1+x^2})\right],
  \label{eq:eos}
\end{align}
where $x = p_F/(m_n c)$ is the Fermi momentum in units of $m_n c$
(Lean: \texttt{Thermodynamics.FermiDirac} + \texttt{Nuclear.BindingEnergy}).

%%============================================================
\section{Derivation of the Maximum Mass}
\label{sec:maxmass}
%%============================================================

\subsection{Structure of the TOV Solution}

The TOV system~\eqref{eq:tov} is integrated from the center
($r = 0$, $P(0) = P_c$) outward until $P(R) = 0$, defining the
stellar radius $R$.  The total gravitational mass is $M = M(R)$.

As the central pressure $P_c$ increases, the total mass $M(P_c)$
increases, reaches a maximum $M_\mathrm{max}$, then decreases.
The maximum corresponds to the onset of instability: configurations
beyond $M_\mathrm{max}$ are dynamically unstable to radial perturbations.

\begin{theorem}[Maximum mass criterion]\label{thm:maxmass}
The maximum stable neutron star mass satisfies
\begin{equation}
  \frac{dM}{dP_c}\bigg|_{M = M_\mathrm{max}} = 0,
  \quad
  \frac{d^2 M}{dP_c^2}\bigg|_{M = M_\mathrm{max}} < 0.
\end{equation}
This is a standard result from stellar stability theory applied to
the TOV system.  The RS derivation follows from J-cost extremization:
$M_\mathrm{max}$ is the saddle point of $\mathcal{F}[P_c]$ in the
TOV cost landscape.
\end{theorem}

\subsection{The Free Neutron Gas Limit}

For the free neutron gas EOS~\eqref{eq:eos} (Oppenheimer-Volkoff
limit), the TOV equation gives

\begin{proposition}[Oppenheimer-Volkoff limit]\label{prop:ov}
With the free neutron gas EOS and no interactions, the maximum mass is
\begin{equation}
  M_\mathrm{OV} = 0.71\,M_\odot,
  \label{eq:mov}
\end{equation}
at a radius $R_\mathrm{OV} = 9.6$ km.  This is a lower bound on the
true maximum mass (nuclear interactions provide additional repulsive
pressure).
\end{proposition}

\subsection{RS Correction from Nuclear Interactions}

Real neutron stars have nuclear interactions that provide additional
pressure beyond the free neutron gas.  In RS, the neutron-neutron
interaction energy is derived from the $\varphi$-ladder:

\begin{definition}[RS nuclear interaction pressure]
\begin{equation}
  P_\mathrm{RS}^\mathrm{int} = P_\mathrm{free}
  \left[1 + \sum_{k \geq 1} c_k \left(\frac{\rho}{\rho_0}\right)^{k/3}
  \varphi^{-k}\right],
  \label{eq:prs}
\end{equation}
where $\rho_0 = 2.3\times10^{14}$ g/cm$^3$ is nuclear saturation density
and $c_k$ are RS-derived coefficients from the nucleon binding energy
curve (Lean: \texttt{Nuclear.BindingEnergyCurve}).
\end{definition}

The leading correction ($k = 1$) gives a stiffening of the EOS that
increases $M_\mathrm{max}$ above~\eqref{eq:mov}.

\begin{proposition}[RS maximum mass estimate]\label{prop:rs_mass}
Including the leading RS nuclear interaction correction,
\begin{equation}
  M_\mathrm{TOV}^\mathrm{RS} \approx 2.0\text{--}2.5\,M_\odot,
  \label{eq:mmax}
\end{equation}
consistent with observations of PSR~J0740+6620 ($2.08\pm0.07\,M_\odot$)
and PSR~J0952-0607 ($2.35\pm0.17\,M_\odot$).

\textbf{HYPOTHESIS} (falsifier: measurement of $M_\mathrm{NS} > 2.8\,M_\odot$
for a neutron star, which would require a significantly stiffer EOS
than the RS $\varphi$-ladder prediction).
\end{proposition}

%%============================================================
\section{Comparison with the Chandrasekhar Limit}
\label{sec:chandrasekhar}
%%============================================================

\begin{center}
\begin{tabular}{lll}
\toprule
Property & White Dwarf (Chandrasekhar) & Neutron Star (TOV) \\
\midrule
Pressure source & Electron degeneracy & Neutron degeneracy \\
Particle mass & $m_e$ & $m_n \approx 1836\,m_e$ \\
Gravity treatment & Newtonian + SR & Full GR required \\
Max mass & $M_\mathrm{Ch} = 1.44\,M_\odot$ & $M_\mathrm{TOV} \approx 2$--$3\,M_\odot$ \\
RS derivation & $\varphi$-ladder electron mass & $\varphi$-ladder neutron mass \\
Lean coverage & \texttt{RS\_Chandrasekhar\_Mass.tex} & This paper \\
\bottomrule
\end{tabular}
\end{center}

The Chandrasekhar mass $M_\mathrm{Ch} = (\hbar c/G)^{3/2}/m_p^2 \times
5.87$ is derived in RS from the proton mass on the $\varphi$-ladder
combined with the electron Compton wavelength (paper:
\texttt{RS\_Chandrasekhar\_Mass.tex}).  The TOV limit involves the same
GR structure but with $m_n$ replacing $m_e/m_p$ as the relevant mass
scale.

%%============================================================
\section{Numerical Results}
\label{sec:numerical}
%%============================================================

\begin{center}
\begin{tabular}{llll}
\toprule
Quantity & RS Prediction & Observation & Status \\
\midrule
Oppenheimer-Volkoff limit & $0.71\,M_\odot$ (free n) & Lower bound & \checkmark \\
PSR J0740+6620 mass & $< M_\mathrm{TOV}^\mathrm{RS}$ & $2.08\pm0.07\,M_\odot$ & Consistent \\
PSR J0952-0607 mass & $< M_\mathrm{TOV}^\mathrm{RS}$ & $2.35\pm0.17\,M_\odot$ & Consistent \\
RS maximum mass & $2.0$--$2.5\,M_\odot$ & $\gtrsim 2.0\,M_\odot$ & HYPOTHESIS \\
Typical NS radius & $10$--$12$ km (from EOS) & $11$--$13$ km (NICER) & Consistent \\
\bottomrule
\end{tabular}
\end{center}

%%============================================================
\section{Lean Formalization Status}
\label{sec:lean}
%%============================================================

\begin{center}
\begin{tabular}{llll}
\toprule
Claim & Lean theorem & Module & Status \\
\midrule
EFE from RS action & \texttt{efe\_from\_mp} & \texttt{Relativity.Dynamics.EFEEmergence} & Proved \\
Schwarzschild metric & \texttt{StaticSpherical} & \texttt{Relativity.StaticSpherical} & Proved \\
Neutron mass on $\varphi$-ladder & \texttt{predict\_mass} & \texttt{Masses.MassLaw} & Proved \\
Fermi-Dirac neutron EOS & \texttt{FermiDirac} & \texttt{Thermodynamics.FermiDirac} & Proved \\
Nuclear binding curve & \texttt{BindingEnergyCurve} & \texttt{Nuclear.BindingEnergyCurve} & Proved \\
Valley of stability & \texttt{ValleyOfStability} & \texttt{Nuclear.ValleyOfStability} & Proved \\
\midrule
TOV max mass from RS EOS & --- & --- & HYPOTHESIS$^\dagger$ \\
NS radius $10$--$13$ km & --- & --- & HYPOTHESIS$^\ddagger$ \\
\bottomrule
\end{tabular}
\end{center}

$^\dagger$ \textbf{HYPOTHESIS}: RS maximum mass $2.0$--$2.5\,M_\odot$
from eq.~\eqref{eq:mmax}; full numerical TOV integration with RS EOS pending.
Falsifier: confirmed detection of $M_\mathrm{NS} > 3\,M_\odot$.

$^\ddagger$ \textbf{HYPOTHESIS}: Typical NS radii from RS EOS and TOV.
Falsifier: NICER/GW measurements showing $R_\mathrm{NS} < 9$ km or
$> 15$ km for canonical $1.4\,M_\odot$ NS.

%%============================================================
\section{Conclusion}
\label{sec:conclusion}
%%============================================================

We have derived the TOV equation from the RS framework (from the
EFE emergence + Schwarzschild static spherical symmetry), combined it
with the $\varphi$-ladder neutron EOS, and obtained the RS prediction
for the maximum neutron star mass $M_\mathrm{TOV}^\mathrm{RS} \approx
2.0$--$2.5\,M_\odot$, consistent with current observations.  The
derivation parallels the Chandrasekhar limit for white dwarfs but
requires the full GR treatment.

\begin{thebibliography}{99}

\bibitem{TOV1939}
R.~C.~Tolman, ``Static solutions of Einstein's field equations for
spheres of fluid,'' \textit{Phys.\ Rev.}\ \textbf{55}, 364 (1939).

\bibitem{Oppenheimer1939}
J.~R.~Oppenheimer, G.~M.~Volkoff, ``On massive neutron cores,''
\textit{Phys.\ Rev.}\ \textbf{55}, 374 (1939).

\bibitem{Fonseca2021}
E.~Fonseca \textit{et al.}, ``Refined mass and geometric measurements
of the high-mass PSR J0740+6620,'' \textit{Astrophys.\ J.\ Lett.}\
\textbf{915}, L12 (2021).

\bibitem{Romani2022}
R.~W.~Romani \textit{et al.}, ``PSR J0952-0607: The fastest and heaviest
known galactic neutron star,'' \textit{Astrophys.\ J.\ Lett.}\
\textbf{934}, L17 (2022).

\bibitem{Washburn2026a}
J.~Washburn, ``The Algebra of Reality: A Recognition Science Derivation
of Physical Law,'' \textit{Axioms} \textbf{15}(2), 90 (2026).

\bibitem{Washburn2026b}
J.~Washburn, ``The Chandrasekhar Mass from the $\varphi$-Ladder,''
preprint (2026).

\end{thebibliography}

\end{document}
